\section{Uniform streaming from rational Choi descriptions}\label{sec:uniformchoi66}
The preceding instrument description is not supplied with a Kraus
factorization over the Gaussian rationals. Computing such a
factorization is unnecessary for classical numerical trajectories. We now allow the
matrix dimension and the entire rational instrument to be input to one
program, rather than fixed constants of that program.

\subsection{Description and exact update}
Fix an encoding of finite lists of signed binary integers. A valid
description consists of $d\ge2$, a positive semidefinite initial matrix with Gaussian-integer entries $P_0$ with trace $z_0>0$, and a finite command list. For each
command $a$ it gives a positive integer $D_a$ and positive semidefinite matrices with Gaussian-integer entries $C_{a,y}$ of order $d^2$ such that
\begin{equation}\label{eq:rationalchoi66}
           \sum_y\operatorname{Tr}_d C_{a,y}=D_a I_d.
\end{equation}
A zero outcome is simply a zero Choi matrix. The associated instrument
has input-first, unnormalized Choi blocks $C_{a,y}/D_a$ and initial state $P_0/z_0$.
Let $S$ be the total encoded length, including dimensions, delimiters,
command identifiers and all matrix entries. The entire description is
retained and charged. In particular the longest identifier is at most
$S$ bits and bounds the parser's temporary token store; a free external
read-only description would define a different space model. For a stored numerator $P$ define
\begin{equation}\label{eq:choiaction66}
 (S_{a,y})_{\alpha\beta}
       =\sum_{i,j=1}^d(C_{a,y})_{i\alpha,j\beta}P_{ij}.
\end{equation}
The indices in this formula are important: no conjugation or
transposition of $P_{ij}$ is inserted.

\begin{lemma}[Integral Choi trajectory update]\label{lem:choitrajectory66}
In the input-first, unnormalized convention
$J(\Phi)=\sum_{i,j}|i\rangle\langle j|\otimes\Phi(|i\rangle\langle j|)$,
for positive semidefinite $P$ with Gaussian-integer entries the matrices in
\eqref{eq:choiaction66} are positive semidefinite matrices with Gaussian-integer entries and
\[
                 \sum_y\tr S_{a,y}=D_a\tr P.
\]
The conditional outcome probabilities are
$\tr S_{a,y}/(D_a\tr P)$, and the next normalized state is
$S_{a,y}/\tr S_{a,y}$ on a positive-mass outcome.
\end{lemma}
\begin{proof}
The Choi convention gives $S_{a,y}=D_a\Phi_{a,y}(P)$, so positivity
follows from complete positivity, and integrality follows from the
entry formula. Taking traces and using \eqref{eq:rationalchoi66}
gives the identity. The scalar factors $D_a$ and any accumulated
previous factors cancel from both probabilities and normalized states.
A positive matrix of trace zero is zero and is never sampled.
\end{proof}

Put $D_* =\max_a D_a$ and, using logarithms to base two in this section,
\begin{align}
 H&=\lceil\log_2(D_*+1)\rceil,&
 \ell_0&=\lceil\log_2(z_0+1)\rceil,\nonumber\\
 Q&=\ell_0+NH,&
 b&=L+\lceil\log_2(6d^2(N+1))\rceil,\qquad
 u=\min\{Q,b\}.\label{eq:uniformparameters66}
\end{align}
The description, once validated, may be kept read-only; charging its
entire length $S$ is included in the bound below.

\begin{lemma}[Reduced Schur entries and bit length]\label{lem:schurminors67}
Let $A$ be a Hermitian matrix of order $q$ with Gaussian-integer entries,
whose real and imaginary parts have absolute value at most $2^h$.
In exact Schur elimination, let $I$ be the ordered indices of the
positive pivots already eliminated. Every remaining entry is
\begin{equation}\label{eq:borderminor67}
 S_{ij}=\frac{\det A[I\mathbin{\|}i,I\mathbin{\|}j]}{\det A[I,I]},
\end{equation}
where $I\mathbin{\|}i$ appends $i$ in that order, and the empty
principal determinant is one. Reduced numerators and denominators
have $O(q(h+\log(q+1)))$ bits. Hermitian positive semidefiniteness is
therefore decidable by this elimination in polynomial bit time and
space in the explicit input length. No optimized polynomial exponent
is asserted.
\end{lemma}
\begin{proof}
When the leading pivot $t$ is positive, block congruence reduces
\[
 \begin{pmatrix}t&b^*\\ b&C\end{pmatrix}\ge0
 \quad\Longleftrightarrow\quad C-bb^*/t\ge0.
\]
A negative pivot rejects. At a zero diagonal, each $2\times2$ principal
minor forces the corresponding row and column to vanish; otherwise
reject, and if they vanish delete that row and column. Positive pivots
make $A[I,I]$ positive definite. The bordered determinant identity
\[
 \det\begin{pmatrix}A[I,I]&A[I,j]\\A[i,I]&A_{ij}\end{pmatrix}
 =\det A[I,I]\,(A_{ij}-A[i,I]A[I,I]^{-1}A[I,j])
\]
proves \eqref{eq:borderminor67}, including after zero rows are skipped.
The specified ordering prevents a spurious permutation sign.

A minor of order $k\le q$ has determinant modulus at most
$k^{k/2}(\sqrt2\,2^h)^k$, by the product of its row norms.
Its real and imaginary parts are integers with
$O(k(h+\log(k+1)))$ bits. The principal denominator is a positive
integer. Cancelling common factors cannot enlarge these lengths.
Each arithmetic operation on reduced Gaussian-rational entries thus
has polynomial operand size. There are polynomially many operations
and at most $q^2$ stored entries; ordinary multiplication, division and
the Euclidean algorithm give the stated bit bounds. For rational
rather than integral input, first clear a common positive denominator.
Its bit length is at most the sum of the encoded denominator lengths,
so this preliminary step also has polynomial cost.
\end{proof}

\begin{theorem}[Variable-description numerical simulation]\label{thm:uniformchoi66}
There is a single classical randomized program with the following
property. Given any valid description above and $N,L\ge2$, it consumes
an adaptive command stream once, emits each outcome immediately, and
retains a legal numerical density matrix. Against every common
classical adaptive controller, its subnormalized history error at
every prefix, and at every common public stopping time bounded by
$N$, is at most $2^{-L}$.

Description validation has polynomial bit time and polynomial working
space in $S$. After validation, the maximum writable space, including
temporary arithmetic and the stored description, is
\begin{equation}\label{eq:uniformspace66}
 O\bigl(S+d^2[H+\log(d+1)+u]+\log(N+1)\bigr),
\end{equation}
with an absolute implicit constant for the fixed machine encoding.
The total peak space is the maximum of this bound and the stated
polynomial preprocessing bound. Expected bit time per command is
$\operatorname{poly}(S)(1+u)^2$ and expected fresh fair-bit consumption
is $O(H+u+\log(d+1))$. The space bound holds on every rejection trial;
running time is not bounded on every such execution. The theorem
requires neither a spectral gap nor a Kraus factorization.
\end{theorem}
\begin{proof}
In grid mode use $B=2^b$ and retain the integer numerator of
$R_B(\rho)$ from Lemma~\ref{lem:positiveround65}; its trace is always
$B+2d^2$. At a command compute the branch integers in
\eqref{eq:choiaction66}, sample an outcome from their exact integer
traces, and repair the chosen normalized branch back to the same
state denominator. The sampler uses fresh binary words of length
$\lceil\log_2(D_a\tr P)\rceil$, rejecting words outside the required
integer interval. Its acceptance probability is greater than one half,
except for the power-of-two case where it is one. Rejected words are
discarded, not recorded. Branch weights can be computed sequentially;
a second scan finds and recomputes the selected branch. Thus the
number of retained matrices does not grow with the number of outcomes.

The proof of Theorem~\ref{thm:instrumentstream65} uses positivity and
\eqref{eq:rationalchoi66}, not rational Kraus entries. It therefore
applies verbatim at the level of its stated block recursion: the
initial error is at most $6d^2/B$, and each subsequent use adds at
most $6d^2/B$, weighted by total simulated branch mass. It gives
\[
       \sum_h\|A_h-\widehat A_h\|_1
                \le(t+1)6d^2/B\le2^{-L}\quad(t\le N).
\]
The outcome probabilities used here are computed from the approximate
stored state, not from the unknown exact state. The block recursion
accounts for that difference. Freezing halted branches proves the same
stopped bound as in Proposition~\ref{prop:posterior66}.

In exact mode retain the selected unnormalized matrix without repair.
This numerator represents a path block up to omitted common scalar
denominators. Those factors cancel in every conditional branch
probability and normalized output; they are not reconstructed or
stored merely to restore a redundant normalization.
Positivity and the trace identity imply
\[
           1\le\tr P_t\le z_0D_*^t<2^{\ell_0+tH}.
\]
Every entry modulus is bounded by the trace. This mode has zero
approximation error. Choose it when $Q\le b$, and choose grid mode
otherwise. The switch is an order bound, not an exact minimization of
small-instance memory.

For space accounting, each Choi numerator entry has modulus at most
$D_a$, by positivity and \eqref{eq:rationalchoi66}. Each summand in
\eqref{eq:choiaction66} is a product of integers of lengths at most
$H$ and $u+O(\log(d+1))$. An accumulator over $d^2$ summands needs
at most $H+u+O(\log(d+1))$ bits. A constant number of $d\times d$
matrices, traces, indices, division buffers and one rejection word
suffice. Initial grid rounding reads the initial entries sequentially;
its extra operand length $\ell_0$ is at most the charged input length
$S$. Counters and address arithmetic give \eqref{eq:uniformspace66}.
Classical multiplication and long division on $w$-bit integers take
$O(w^2)$ bit operations and $O(w)$ scratch. Scanning the explicit
outcome list a fixed number of times gives the claimed polynomial
factor in $S$; powers of $H$ are absorbed there. The sampling bound
follows from the rejection interval length and the expected number of
trials, which is below two. A single nonzero branch requires no random
bits.

Validation uses the reduced-rational Schur procedure of
Lemma~\ref{lem:schurminors67}, not the unreduced fixed-dimensional
scaled recursion. Its ratios-of-minors formula proves polynomial
operand lengths. Hermitian symmetry, the exact partial trace
identities and positive initial trace are checked separately. No
pivot permutation is needed: a zero leading diagonal of a positive
semidefinite matrix forces its entire row and column to vanish.
This establishes a separate polynomial preprocessing bound, not a
validator matching the subsequent streaming-space bound.

Parameters are parsed before commands. The precision $L$ may be read
with saturation at $Q$: once that threshold is reached, exact mode is
selected and its exact value is unnecessary. The time to read all
parameter digits is charged as preprocessing in their actual length;
no equally long precision integer is retained. A binary counter checks
the command count. The numerical output consists of signed binary
integer entries and their common denominator; writing it requires no
additional matrix-sized precision store.
\end{proof}

\begin{corollary}[Known joint systems and two error budgets]\label{cor:jointsimulation66}
Suppose the entire joint system to be simulated, including a specified
reference and all specified controller operations, has a rational Choi
description on dimension $d_{\rm joint}$. Theorem~\ref{thm:uniformchoi66}
applies with that dimension and the full description length charged.
If its instruments were first compiled by
Theorem~\ref{thm:choiround66}, the reference-stable description error
and the numerical history error add. In particular, allocating
$\varepsilon/2$ to each gives total history error at most $\varepsilon$.
\end{corollary}
\begin{proof}
Use the joint Choi maps and joint initial state as input to the uniform
program. Theorem~\ref{thm:adaptivediamond66} bounds the first
replacement, and Theorem~\ref{thm:uniformchoi66} bounds the numerical
simulation of the replaced experiment. The triangle inequality gives
the sum. An unspecified quantum tester remains covered by the
\emph{instrument-description} comparison only; a classical numerical
simulation of it is not thereby supplied at no cost.
\end{proof}

Theorem~\ref{thm:uniformchoi66} is an upper bound with variable data.
The lower bounds of Theorem~\ref{thm:matrixspace65} retain their own
fixed-program quantifiers, numerical matrix interface, and full-action
spectral and cap hypotheses. No dimension-uniform spectral constant,
no sharp lower bound for arbitrary noisy transcript-only interfaces,
and no classification of disturbing instruments follows from the
upper construction. Conversely, a rational Choi input is enough to
execute the construction even when a Gaussian-rational Kraus list is
not supplied. This is a representation-level enlargement of the
constructive theorem, not an alteration of the inherited lower model.
